\documentclass[a4paper]{article}

\usepackage[justification=centerlast]{caption}
\usepackage[utf8]{inputenc}   % LaTeX, comprends les accents !
\usepackage[T1]{fontenc}      % Police contenant les caractères français
\usepackage[french]{babel}
\usepackage{graphicx}  % pour inclure des images
\usepackage{xcolor}    % pour utiliser des couleurs
\usepackage{amsmath} % Pour les mathématiques
\usepackage{minted} % Pour les blocs de code (assurez-vous que Pygments est installé)
\usepackage{geometry} % Pour définir les marges
\usepackage{listings} % Pour la coloration syntaxique du code
\usepackage{hyperref} % Pour les liens cliquables
\usepackage{subcaption}
\usepackage{float} %pour utiliser [H]

% Configuration de listings (pour code et autres)
\lstset{
    language=Python, % Langage du code
    basicstyle=\ttfamily\small, % Style de base
    keywordstyle=\color{blue}, % Style des mots-clés
    commentstyle=\color{gray}, % Style des commentaires
    stringstyle=\color{red}, % Style des chaînes de caractères
    numbers=left, % Numérotation des lignes à gauche
    numberstyle=\tiny\color{gray}, % Style des numéros de ligne
    stepnumber=1, % Incrément des numéros de ligne
    numbersep=5pt, % Espacement des numéros de ligne
    backgroundcolor=\color{white}, % Couleur de fond
    frame=single, % Cadre autour du code
    breaklines=true, % Casse les lignes trop longues
    breakatwhitespace=false, % Ne casse pas les lignes aux espaces
    tabsize=4, % Taille des tabulations
    captionpos=b, % Position de la légende
    showspaces=false, % Ne montre pas les espaces
    showstringspaces=false, % Ne montre pas les espaces dans les chaînes
    showtabs=false, % Ne montre pas les tabulations
}

% Définir les marges avec le package geometry
\geometry{
    a4paper, % Format du papier
    top=2cm, % Marge supérieure
    bottom=2.5cm, % Marge inférieure
    left=2.5cm, % Marge de gauche
    right=2.5cm % Marge de droite
}

\title{Projet de programmation L3\\[0.5cm]Développement de méthodes pour un logiciel d'extraction de courbes temps-activité vasculaire en tomographie d'émission de positons (TEP)}

\author{Alia CHEBBI, Camille GUIGUEN, DIEVART Colin, Imane LAZIZI, Kanzy YOUSSEF\\[0.5cm]
    \emph{Encadrant :}Dr.Florentin KUCHARCZAK\\
    \emph{En lien avec :} CHU de Montpellier\\[0.5cm]
    \emph{L3 informatique}\\Faculté des Sciences\\Université de Montpellier. }
\date{2024-2025}

\begin{document}

\begin{titlepage}
    \begin{minipage}{0.18\textwidth}
        \includegraphics[width=\linewidth]{UM_logo.jpg} % Remplacez par le chemin de votre logo
    \end{minipage}
    \hfill
    \begin{minipage}{0.3\textwidth}
        \centering
        \includegraphics[width=\linewidth]{CHU_logo.jpg} % Remplacez par le chemin de votre logo
    \end{minipage}
    \hfill
    \begin{minipage}{0.3\textwidth}
        \raggedleft
        \includegraphics[width=\linewidth]{fds_logo.jpg} % Remplacez par le chemin de votre logo
    \end{minipage}

    \vspace{2cm}

    \centering
    {\scshape\LARGE Université de Montpellier \par}
    \vspace{1cm}
    {\scshape\Large Projet de programmation L3\par}
    \vspace{1.5cm}
    {\huge\bfseries Développement de méthodes pour un logiciel d'extraction de courbes temps-activité vasculaire en tomographie d'émission de positons (TEP)\par}
    \vspace{2cm}
    {\Large\itshape Alia CHEBBI, Camille GUIGUEN, DIEVART Colin, Imane LAZIZI, Kanzy YOUSSEF\par}
    \vfill
    \emph{Encadrant :} Dr.Florentin KUCHARCZAK\par
    \emph{En lien avec :} CHU de Montpellier et l'Université de Montpellier \par
    \emph{Cadre :} UE HAI606I  Projet de Programmation 2\par
    \vfill
    {\large 2024-2025\par}
\end{titlepage}

\newpage
\dominitoc  
\tableofcontents
\newpage
\listoffigures
\newpage


%%%--------------------------------------------INTRODUCTION-------------------------------------------------%%%

\section{Introduction} 
    Dans le cadre de ce projet de programmation, nous avons collaboré avec le service de médecine nucléaire du CHU de Montpellier pour développer des nouvelles méthodes à implémenter dans un logiciel d’analyses médicales.\\

    \noindent Ce logiciel vise à traiter des données dynamiques\footnote{Les données dynamiques ou transactionnelles (ou encore "dyata") sont des informations mises à jour périodiquement, ce qui signifie qu'elles changent de manière asynchrone au fil du temps.\\} obtenues par Tomographie par Émission de Positons (TEP), une technique d'imagerie médicale permettant de visualiser et de mesurer l'activité métabolique des organes et des tissus du corps. On injecte au patient un médicament radio-pharmaceutique contenant des traceurs renvoyant des "signaux"\footnote{Après injection du Fluorodeoxyglucose (le traceur utilisé dans ce projet), la désintégration de ce dernier et après une suite d'évènements physiques, nous aboutissons à l'émission de deux photons dans une direction commune mais de sens opposé. Les coïncidences sont captées par des anneaux de détecteur et ensuite reconstruite à l'aide d'algorithmes complexes.} pour créer des images tridimensionnelles représentant la distribution du radiotraceur dans les tissus.
    Les données des images reconstruites sont enregistrées dans des fichiers NifTI.\\
    
    \noindent Ce logiciel, utilisé dans le cadre de recherches dans le service de médecine nucléaire, nécessite encore cependant l'utilisation de briques non automatisées et faisant appel à des logiciels externes pour : 
    \begin{enumerate}
        \item Le recalage du cerveau dans un espace normalisé dans lequel les correspondances anatomiques des voxels est identifiée.
        \item L'extraction de la fonction d'entrée vasculaire.
    \end{enumerate}
    Ces deux étapes sont pour l'instant effectuées à l'aide d'un logiciel propriétaire (pmod). \\
    
    \noindent \textbf{L'objectif final de ce projet est de remplacer la brique d'extraction manuelle de la fonction d'entrée vasculaire par une version automatisée.}\\ 
    Plus précisément ces méthodes doivent estimer des courbes d'entrée vasculaire Temps-Activité (TAC) dans la région de la carotide. Entre autres, analyser la concentration du médicament dans la carotide grâce aux données des fichiers Nifti qui nous ont été fournis.
    Une fois ces deux briques implémentées et validées en comparaison avec le gold standard, ce logiciel a vocation à être publié (dans une revue comme \href{}{Neuroinformatics} par exemple) et ensuite rendue accessible via un Git.\\

    \noindent Ce projet en particulier a attiré notre attention car il se distingue des projets de programmation standards comme le 
    développement de jeux. Il nous permet de découvrir l’application de l’informatique dans un domaine nouveau qui nous 
    intéresse, la médecine nucléaire. De plus, cela nous donne un bon exemple de l’utilité parfois vitale de la programmation.\\

    \vspace{0.3cm}

\noindent\underline{\textbf{CAHIER DES CHARGES :}}

    \vspace{0.2cm}
    \noindent\textbf{Objectif final :}\\
    Implémentation de méthode(s) provenant d’un algorithme que nous jugeons optimal par rapport aux pistes que nous avons étudiées.
    \begin{itemize}
        \item \textbf{En entrée :} un fichier au format NifTI contenant des données patients déjà recadrées de la zone du cerveau et comprenant la zone carotidienne.
        \item \textbf{En sortie :} une ou plusieurs courbes temps-activité (TAC) caractéristiques de la zone carotidienne.
    \end{itemize}

    \vspace{0.2cm}
    \noindent\textbf{Problématiques :}
    \begin{itemize}
        \item Difficultés scientifiques pour des novices du domaine de la médecine nucléaire.
        \item Quelle méthode optimale choisir et où trouver de la bonne documentation.
        \item Comment extraire des clusters uniquement de la zone carotidienne des images TEP.
        \item La gestion des patients avec des anomalies (morphologique ou dans la présence d'artefacts comme un cathéter dans le champ de vue).
    \end{itemize}
    
    \vspace{0.2cm}
    \noindent\textbf{Critères de succès pour le code final:}
    \begin{itemize}
        \item Génération de TAC en sortie.
        \item Automatisation complète.
        \item Fonctionnement avec tout les patients fourni.
        \item Meilleures performances que la méthode déjà existante.\\
    \end{itemize}
    
    
    \noindent\textbf{Acteurs du projet :}\\
    \textit{Encadrant} : Dr. Florentin KUCHARCZAK\\
    \textit{Professeur référente} : Dr. Anne-Elisabeth BAERT\\
    
    \noindent\textbf{Livrables attendus :}
    \begin{itemize}
        \item Des "rapports d'avancements" tout au long du projet (disponible en annexe).
        \item Un rapport final de projet, d'une vingtaine de pages, au format Latex.
        \item Un fichier Python consultable, organisé et commenté validant les objectifs fixés.
        \item Un espace GitHub structuré comprenant toutes les pistes étudiées.
        \item Une soutenance orale finale, en mai 2025.\\
    \end{itemize}


%%%--------------------------------------------mainPart-----------------------------------------------%%%
\newpage
\section{Algorithmes et chemins de pensée}
    \subsection{Trouver une méthode générale}
            \noindent
            \label{sec:intro}
            Un des plus grands défis de ce projet a été, pour les informaticiens en herbe que nous sommes, de comprendre les objectifs du projet d’un point de vue scientifique. Nous avons pris appui sur les articles de référence que notre encadrant nous a transmis (voir \hyperref[bibliographie]{\textbf{la bibliographie}}), nous nous sommes familiarisé avec le logiciel \textbf{ITK-SNAP} (voir \hyperref[sec:Outils]{\textbf{la partie 3}} pour plus de détails).\\
            
            \noindent
            Afin d’aborder la problématique des images TEP dynamiques en 4D, nous avons d’abord exploré manuellement les fichiers d’acquisition. Grâce à une première série de scripts Python, nous avons visualisé les images en coupe axiale, coronale et sagittale du fichier \texttt{w\_PET\_4D.nii} du patient P01. Pour mieux appréhender la dimension temporelle, nous avons extrait les valeurs d’intensité pour chaque voxel $(x, y, z, t)$, trié les intensités supérieures à un seuil, et sauvegardé ces données volumineuses dans un fichier brut . Ce fichier nous a servi de base de comparaison pour valider les courbes temps-activité (TAC) obtenues ensuite.\\
            
            \noindent
            Nous avons également représenté dans nos courbes (TAC) les valeurs du modèle de Feng (souvent en pointillés), un modèle théorique standard qui ajuste nos données, afin de permettre une comparaison avec les données extraites dans la région d'intêret (ROI\footnote{R.O.I : Region Of Interest}) . Un lien vers un article expliquant ce modèle en détail \href{https://www.sciencedirect.com/science/article/abs/pii/002071019390049C?via\%3Dihub}{\textbf{ici}}.\\
            
            \noindent
            Après cette phase d'analyse bibliographique, et la lecture attentive des articles comme : \textbf{\textit{\href{https://pubmed.ncbi.nlm.nih.gov/37947880/}{An update on the use of image-derived input functions for human PET studies: new hopes or old illusions?}}}, où plusieurs méthodes d’estimation de la fonction d’entrée nous ont été présentées :

            \begin{enumerate}
                \item \textbf{SIME (Simultaneous Estimation)} : une estimation en même temps que les paramètres cinétiques et la fonction d'entrée artérielle à partir des courbes temps-activité (TAC) issues de l'image. Cette méthode a une complexité computationnelle beaucoup trop élevée pour nos machines.
                \item \textbf{PBIF (Population-Based Input Function)} : Au lieu de mesurer directement l'activité du traceur dans le sang, cette approche utilise une courbe moyenne dérivée d'une population de référence, ajustée au sujet grâce à des échantillons de sang artériel ponctuels. Ce qui n'est pas envisageable en routine clinique car trop invasif. En effet, le prélèvement artériel est plus délicat et douloureux que le prélèvement veineux.
                \item \textbf{IDIF (Image-Derived Input Function)} : une méthode cette fois non invasive basée sur la segmentation de structures vasculaires dans l’image, et le \textit{clustering} des données dynamiques. C’est cette dernière que nous avons choisie d’explorer et d’implémenter comme méthode principale.\\
            \end{enumerate}
            
            \noindent
            Parmi les différentes variantes testées de l’IDIF, nous avons estimé que l’approche basée sur le \textit{clustering} était la plus adaptée à notre cas. Cet algorithme consiste à regrouper les zones d’intérêt (ROIs) de l’image TEP en fonction de similarités temporelles et spatiales, selon des critères prédéfinis.\\
            
            Le processus d'extraction de l'IDIF via \textit{clustering} se déroule en trois grandes étapes :
            
            \begin{enumerate}
                \item \textbf{Clustering} des voxels pour regrouper les régions dynamiquement similaires ;
                \item \textbf{Segmentation} des vaisseaux d’intérêt à partir de ces clusters ;
                \item \textbf{Correction}, notamment des effets de volume partiel (PVE), pour obtenir une estimation plus fidèle de la courbe d’entrée. (Modéle de Feng)\\
            \end{enumerate}
            
            \noindent
            Nous avons également comparé différentes approches d’affichage des TAC, en analysant les courbes extraites pour un voxel spécifique et pour un cube de voxels avoisinants, afin d’évaluer la stabilité et la représentativité du signal. En parallèle, nous avons développé une méthode d’extraction semi-automatique des valeurs de la carotide, d’abord sans bibliothèque externe (dans le fichier \href{run:./Pistes_et_Tests/autres_tests/test.py}{\texttt{test.py}}), puis en intégrant des outils comme \textbf{Nilearn} pour l’automatisation de la segmentation, l’extraction de masques, et le traitement spatio-temporel. Nous avons également exploré le module de parcellisation de Nilearn, ouvrant des perspectives intéressantes pour une segmentation plus robuste et intelligente. \\
\newpage
        \item \textbf{Approche intuitive et limitations pour localiser les carotides}
    
            Avant de passer aux méthodes de clustering et de segmentation, nous avons exploré des approches simples pour identifier au moins un voxel (ou un voxel + ses voisins) appartenant aux artères carotides, afin d’extraire leurs TAC\footnote{T.A.C : Time-Activity Curve} et analyser les graphes obtenus.
    
            En nous basant sur nos connaissances anatomiques, nous avons restreint les volumes 4D (x, y, z, t) au plan transversal (cf figure 2), en coupant les valeurs les plus basses de l’axe z (cf figure 1), puisque les carotides sont situées sous le cerveau, au niveau du cou.
    
            \vspace{0.5cm}

            \begin{figure}[H]
                \centering
                \includegraphics[width=0.9\linewidth]{p11 frontal z=8.png}
                \caption{Coupe frontale}
                \label{fig:frontal}
            \end{figure}
            
            \noindent
            On distingue les carotides sous la ligne z=8 (zones rouges) (cf figure 1).
            
            \vspace{0.5cm}
            
            \begin{figure}[H]
                \centering
                \includegraphics[width=0.9\linewidth]{p11 transversal z=8.png}
                \caption{Coupe transversale}
                \label{fig:transversale}
            \end{figure}
            
            \noindent
            On observe des zones de forte intensité (rouge), indiquant une activité métabolique élevée (cf figure 2).


            Comme cette zone contenait peu de tissus (principalement vasculaire et musculaire), nous avons supposé qu$’$un simple seuil d’intensité permettrait de distinguer les régions d’intérêt, ce qui nous a conduit à créer un masque binaire. Cette approche intuitive a été testée dans les scripts \href{run:./Pistes_et_Tests/autres_tests/pistes_tests/masking_thresholding_each_patient.py}{\texttt{masking\_thresholding.py}} et \href{run:./Pistes_et_Tests/autres_tests/pistes_tests/standard_masking.py}{\texttt{standard\_masking.py}} (cf figure 3). Nous nous concentrons surtout sur ce dernier, qui a donné des résultats pertinents.


            Nous avons choisi un patient de référence sans anomalies visibles (vérifié sous ITK\-SNAP, patient w\_PET\_4D ou P00.nii), afin de définir des paramètres fixes et fiables de seuil (par ex. 23) et de découpe en z. Ce masque binaire a ensuite été appliqué aux autres patients, après recalage et rééchantillonnage, puis en sélectionnant les voxels avec une intensité >0. Enfin, nous avons identifié le voxel à intensité maximale et calculé l’intensité moyenne avec ses voisins (cf figure 4) .
            \vspace{0.5cm}

            \begin{figure}[H]
                \centering
                \includegraphics[width=0.9\linewidth]{exemple masque standard.png}
                \caption{Résultat de masque binaire}
                \label{fig:masque_binaire}
            \end{figure}
            
            
            \vspace{0.5cm}
            
            \begin{figure}[H]
                \centering
                \includegraphics[width=0.9\linewidth]{P01_TAC_fit.png}
                \caption{Exemple de la courbe TAC}
                \label{fig:exemple_TAC}
            \end{figure}
            
         
            Cette approche n’était qu’un test exploratoire, réalisé pour mieux comprendre la problématique tout en évitant les coûts computationnels et les surcharges mémoire.
            
            \noindent
            Elle a également mis en évidence les limites d’un découpage (crop) et de seuils statiques, qui peuvent introduire des erreurs futures. Par exemple, ce code échouerait sur des scans contenant l’ensemble du torse, ou sur des scans non alignés avec le patient de référence (comme P00.nii).
            
            \noindent
            Un autre inconvénient est qu’il nécessite de sélectionner manuellement les patients et les paramètres, ce qui le rend peu pratique et fragile.
            
            \noindent
            Pourtant, ces petites failles ont permis d’identifier des défis importants, comme repérer des «patients irréguliers» (P03, P04, P09), où les zones les plus lumineuses correspondaient en réalité au cathéter à l’extérieur du cerveau (cf figure 5 et figure 6).
            
            \begin{figure}[H]
                \centering
                \begin{minipage}[t]{0.48\linewidth}
                    \centering
                    \includegraphics[width=0.7\linewidth]{P03.png}
                    \caption{Exemple sur patient 3}
                    \label{fig:patient3}
                \end{minipage}
                \hfill
                \begin{minipage}[t]{0.48\linewidth}
                    \centering
                    \includegraphics[width=0.7\linewidth]{P09.png}
                    \caption{Exemple sur patient 9}
                    \label{fig:patient9}
                \end{minipage}
            \end{figure}

            \noindent
            Pour y remédier, nous avons introduit des masques spécifiques (par ex. \texttt{brain\_mask}, \texttt{lower\_brain\_mask} (cf figure 11 et figure 10)) afin d’éviter de détecter le cathéter au lieu des carotides.

            \noindent
            Enfin, cette phase a montré les limites de l’approche naïve par seuillage et a motivé l’utilisation de méthodes plus robustes comme le clustering.
            
 
        \end{itemize}
    \subsection{Les différentes méthodes de clustering}
        
            Dans cette section, nous explorons trois méthodes de clustering appliquées à la segmentation de la carotide : K-Means, Mean Shift, et Ward.\\
            \noindent
            Chaque méthode est analysée en termes d'objectifs, de mise en œuvre, d'avantages et d'inconvénients. Nous verrons à la fin de cette section quelle à été la méthode retenue et pourquoi.\\
    
            \noindent Notre objectif est d'identifier et d'isoler automatiquement la zone de la carotide puis de détecter les clusters les plus pertinents.\\
            \noindent Pour les clusters, la bibliothèque NiLearn propose une méthode \href{https://nilearn.github.io/dev/modules/generated/nilearn.regions.Parcellations.html}{\textcolor{blue}{\texttt{Parcellation()}}} très complète.\\
            \noindent Cette méthode consiste à appliquer un des 5 algorithmes différents de clustering (cf Figure 7) à des données d’imagerie cérébrales (fMRI) pour diviser le cerveau en régions distinctes qu’on appelle ‘clusters’ ou ‘parcelles’ : des composantes spatiales du cerveau où l'intensité a des ‘patterns’ temporels similaires ou des points qui sont cohérents ensemble.\\
            
            \begin{figure}[H]
                \centering
                \includegraphics[width=0.9\linewidth]{all_algo_parcellation.png}
                \caption{Résultats des différentes méthodes de la fonction \texttt{Parcellation()}}
                \label{fig:image}
            \end{figure}
            
  %%%-------------------------------------------K-MEANS---------------------------------------------------%%%
\newpage  
        \subsubsection{CLUSTERING K-MEANS  \href{run:/Pistes_et_Tests/piste_kmeans/testscikit.py}{(\underline{\texttt{testscikit.py}}})}}
        Nous allons commencer explorer la piste de l'algorithme K-Means, une méthode de clustering largement utilisée pour regrouper des données en clusters basés sur des similarités. 
        Nous nous concentrons sur la segmentation de la carotide, car les artères carotides sont les principaux vaisseaux sanguins qui fournissent le sang au cerveau. Pour cela il nous faut d'abord réduire la zone d'analyse.\\

        \subsubsection*{Filtrage des données}
        \noindent
        Pour réduire le nombre de données à analyser, on choisit ici de supprimer les 90\% des voxels avec l’intensité la plus basse, ce qui nous laisse les voxels les plus concentrés et donc les plus intéressants à analyser.
        
%NE PAS MODIFIER : indentation moche mais obligatoire pour que le code s'affiche bien                       
            \begin{lstlisting}[language=Python]
# Filtrer les voxels ayant une intensité significative (évite le bruit)
threshold = np.percentile(img_4D, 90)  #Garde les 10% de valeurs les plus hautes
mask = img_4D.max(axis=-1) > threshold
            \end{lstlisting}
            \\
            
            \noindent
            Ici on met dans un tableau appelé \texttt{voxel\_time\_series}, qui pour chaque voxel utilisé va conserver ses 29 valeurs au cours du temps pour appliquer l’algorithme K-means dessus.\\
            
            \begin{lstlisting}[language=Python]
# Extraire les courbes temporelles des voxels sélectionnés
x, y, z = np.where(mask)
voxel_time_series = img_4D[x, y, z, :]
            \end{lstlisting}
            
            \subsubsection*{Algorithme de base}
            Voici le pseudo-code de l'algorithme KMeans :
            
            \begin{itemize}
                \item Extraire l'espace des features de l'image
                \item Initialiser $k$ centres aléatoire
                \item Répéter :
                \begin{itemize}
                    \item Pour chaque point :
                    \begin{itemize}
                        \item calculer la distance à chaque centre
                        \item assigner le point au centre le plus proche
                    \end{itemize}
                    \item Pour chaque cluster :
                    \begin{itemize}
                        \item recalculer le centre comme la moyenne des points assignés
                    \end{itemize}
                    \item calculer le déplacement de chaque centre
                \end{itemize}
                \item Tant que déplacement maximal $>$ Seuil \textbf{et} nombre d’itérations $<$ \texttt{max\_iter}
            \end{itemize}

                
           
            \subsubsection*{Explication de K-Means}
                Maintenant que les données sont filtrées, on applique l'algorithme K-means (méthode des "Plus Proches Voisins", vu dans le module d'Algorithmique).\\
    
                K-Means est un algorithme de clustering non supervisé, c'est-à-dire qu'il permet de regrouper automatiquement des données en groupes (ou "clusters"), sans avoir besoin d'étiquettes (contrairement au classifieur supervisé).
                
                Il est utilisé pour regrouper des objets similaires selon leur forme, intensité, ou comportement, comme ici dans le cas de voxels avec des courbes temporelles similaires.\\
                
                L'algorithme commence par choisir $K$ centres de cluster (appelés centroïdes). Chaque point (dans notre cas : un voxel représenté par sa courbe temporelle) est assigné au cluster dont le centre est le plus proche.
                
                \noindent La distance utilisée est généralement la distance euclidienne, c'est-à-dire :
                \[ d(p, c) = \sqrt{(p_1 - c_1)^2 + (p_2 - c_2)^2 + \cdots + (p_n - c_n)^2} \]
                
                Une fois tous les points assignés à un cluster, on recalcule le centre de chaque cluster. On répète ce procédé jusqu'à ce que les résultats de K-means ne bougent plus et qu’on obtienne donc nos centres. Ici, on choisit 3 centres, afin de maximiser les chances de trouver la carotide, car on a plusieurs clusters possibles : la carotide, les tissus du cerveau, et les vaisseaux secondaires.
            
            \subsubsection*{Résultats à 3 clusters}
                Ici, on remarque que 2 clusters (cf Figure 8: Clusters 0 et 1) ont des coordonnées très proches, ce qui rend inutile un cluster. 
                On décide donc de passer à 2 clusters afin de voir si ce sera plus précis (car avec 2 clusters ils ne se chevaucheront pas).
           
            %figure 2 - Res 3 clusters
            \begin{figure}[H]
                \centering
                \includegraphics[width=0.90\linewidth]{res_3_cluster.png}
                \caption{TAC résultantes de 3 clusters (0, 1 et 2)}
                \label{fig:image1}
            \end{figure}
            
            \subsubsection*{Résultat à 2 clusters}
            On remarque qu'un nouveau cluster avec des coordonnées différentes apparaît (cf Figure 9 : courbe verte, Cluster 1).\\
            \noindent Lorsqu'on regarde sur ITK-SNAP, on peux voir qu'il appartient à la carotide gauche.
            %figure 3 - Res 2 cluster
            \begin{figure}[H]
                    \centering
                    \includegraphics[width=0.90\linewidth]{res_2_cluster.png}
                    \caption{TAC résultantes avec 2 cluster (0 + TAC un nouveau cluster)}
                    \label{fig:enter-label}
                \end{figure}                
                    
            
            \subsubsection*{Conclusion sur K-Means}
            Après avoir peaufiné l'algorithme et effectué des modifications, nous sommes arrivé à des conclusions sur cette méthode et son efficacité.
            
            \subsubsection*{Les avantages}
            L'algorithme est simple et constant, il est adapté pour les données en grand nombre, il trouvera toujours des clusters.
            
            \subsubsection*{Les inconvénients}
            \begin{enumerate}
                \item Automatisation et instabilité\\
                Le nombre de clusters doit être fixé à l'avance, la distance euclidienne n'est pas adaptée aux données spatio-temporelles. Les clusters sont trop imprécis, et les résultats diffèrent trop en fonction du nombre de clusters choisi, ce qui fait que l'algorithme pour ce cas-là n'est pas très fiable.\\
                \item Forme des clusters\\
                Les clusters sont sphériques, ce qui n'est pas vraiment le cas de la carotide, et donc cela réduit les chances de trouver la carotide.\\
            \end{enumerate}
         
            Pour conclure, l'algorithme K-Means nous a permis de travailler sur une méthode différente, plus algorithmique mais pas forcément efficace et précise pour ce genre de cas. La méthode n'est donc finalement pas viable pour cette situation, mais elle nous a aidé à comprendre les méthodes de Clustering pour la suite.
            
    \hspace{1cm}   
    %\hrule
    \newpage

%%%-------------------------------------------MEAN SHIFT----------------------------------------------------%%%

       
        \subsubsection{CLUSTERING MEAN SHIFT \href{run:/Pistes_et_Tests/piste_meanshiftetbrainmask/MeanShift/MeanShift.py}{(\underline{\texttt{MeanShift.py}})}}

        \emph{Le clustering Mean Shift (cf Figure 10) est une technique non paramétrique qui ne nécessite pas de spécifier le nombre de clusters. Robuste aux outliers\footnote{Outliers : points aberrants ou isolés qui diffèrent significativement de la majorité des données}, il ne suppose pas que les clusters sont sphériques, ce qui lui permet de s'adapter à des formes complexes et non linéaires.}
        
        \subsubsection*{Algorithme de base}
        Voici le pseudo-code de l'algorithme Mean Shift :
        
        \begin{itemize}
            \item Extraire l'espace des features de l'image
            \item Tant que points non visités $>$ 0 :
            \begin{itemize}
                \item mean $\gets$ point aléatoire
                \item Répéter :
                \begin{itemize}
                    \item dist $\gets$ distance(mean, tous les points)
                    \item neighbors $\gets$ points dans la bandwidth
                    \item new\_mean $\gets$ moyenne(neighbors)
                    \item Si dist(new\_mean, mean) $<$ Seuil :
                    \begin{itemize}
                        \item assigner les neighbors au cluster
                        \item marquer les neighbors comme visités
                        \item sortir de la boucle
                    \end{itemize}
                    \item mean $\gets$ new\_mean
                \end{itemize}
            \end{itemize}
        \end{itemize}
        
        \begin{figure}[H]
            \centering
            \includegraphics[height=6cm]{mean_shift-exemple.png}
            \hspace{0.5cm}
            \includegraphics[height=6cm]{mean_shift_tutorial-38.png}
            \caption{Illustration du clustering \textit{Mean Shift}}
            \label{fig:meanshift}
        \end{figure}
        
        \footnotetext{Source des images figure 10: \href{https://sbme-tutorials.github.io/2019/cv/notes/6_week6.html#mean-shift-clustering}{sbme-tutorials.github.io}}



        
        \subsubsection*{Pourquoi Mean Shift pour la segmentation carotidienne ?}
            Mean Shift a été choisi car :
            \begin{enumerate}
            	\item Il détecte automatiquement le nombre de clusters.
            	\item Robuste aux formes irrégulières : Les carotides n’ont pas une forme sphérique parfaite et l’intensité locale peut varier. MeanShift s’adapte à des clusters de formes non linéaires, ce qui est un avantage majeur pour capter ces structures vasculaires.

                \item Adapté aux zones denses : MeanShift identifie les maxima locaux de densité. Comme les carotides apparaissent comme des régions intenses et denses dans les images prétraitées, l’algorithme est naturellement dirigé vers ces zones.

                \item Il gère bien les données médicales bruitées et irrégulières
            \end{enumerate}
    
        \subsubsection*{Implémentation : Segmentation carotidienne}
            \subsubsection*{Pré-traitement}
                Pour améliorer le clustering, nous avons restreint l'analyse à une petite région,vue que les artères carotides intracrâniennes représentant moins de 0.2\% du volume cérébral total.

                Anatomiquement, les carotides sont situées :
                \begin{itemize}
                    \item Sous le cerveau (région cervicale)
                    \item Près de la ligne médiane du corps
                    \item Ni trop latérales ni spécifiquement antérieures/postérieures
                \end{itemize}

                D'où l'utilisation de trois masques successifs :
                
              \begin{figure}[H]
                \centering
                \includegraphics[width=0.7\linewidth]{masque du cerveau.png}
                \caption{Masque cérébral}
                \label{fig:masque_cerveau}
            \end{figure}
            \noindent
            \small Pour isoler les voxels cérébraux (cf Figure 11).
            
            \vspace{0.5cm}
            
            \begin{figure}[H]
                \centering
                \includegraphics[width=0.7\linewidth]{partie inférieur du cerveau.png}
                \caption{Masque inférieur}
                \label{fig:masque_inferieur}
            \end{figure}
            \noindent
            \small Restreint à la région inférieure (cf Figure 12).
            
            \vspace{0.5cm}
            
            \begin{figure}[H]
                \centering
                \includegraphics[width=0.7\linewidth]{zone carotidienne.png}
                \caption{Masque carotidien}
                \label{fig:masque_carotidien}
            \end{figure}
            \noindent
            \small Zone d'intérêt carotidienne (cf Figure 13).



            \subsubsection*{Segmentation des artères carotides (cf Figure 14)}
                Le processus comporte deux étapes clés :\\

                \textbf{Traitement par hémisphère (process\_hemisphere)} :
                \begin{itemize}
                    \item Applique le clustering MeanShift pour séparer les structures vasculaires brillantes
                    \item Traite un hémisphère à la fois (gauche/droite du plan sagittal médian)
                    \item Utilise des features 4D : coordonnées spatiales (x,y,z) + intensités
                    \item Estime automatiquement la bandwidth si non spécifiée
                    \item Retourne les deux clusters les plus brillants (carotide interne et externe)
                \end{itemize}
                
                \vspace{1em} 
                \hspace{0em}\textbf{Extraction carotidienne (extract\_carotids)} :
                \begin{itemize}
                    \item Identifie la frame temporelle optimale avec le pic de rehaussement artériel
                    \item Sépare le masque en hémisphères gauche/droit
                    \item Traite chaque hémisphère via process\_hemisphere
                    \item Combine les résultats pour créer un masque carotidien complet
                \end{itemize}

            Cette combinaison de restriction spatiale (via masques anatomiques) et de clustering par intensité (via MeanShift) permet une segmentation robuste des carotides tout en minimisant les faux positifs (cercle de Willis, artères vertébrales (cf Figure 15)).\\
            
        \noindent

    \begin{figure}[H]
        \centering
        \begin{subfigure}[t]{0.45\textwidth}
            \centering
            \includegraphics[width=\linewidth]{4-left_carotid.png}
            \caption{Masque carotide gauche}
            \label{fig:carotide_gauche}
        \end{subfigure}
        \hfill
        \begin{subfigure}[t]{0.45\textwidth}
            \centering
            \includegraphics[width=\linewidth]{4-right_carotid.png}
            \caption{Masque carotide droite}
            \label{fig:carotide_droite}
        \end{subfigure}
    
        \vspace{0.5cm}
    
        \begin{subfigure}[t]{0.7\textwidth}
            \centering
            \includegraphics[width=\linewidth]{detailed_carotid_mask.png}
            \caption{Masque combiné}
            \label{fig:carotide_combine}
        \end{subfigure}
    
        \caption{Visualisation des masques des carotides gauche, droite et combinée.}
        \label{fig:carotides}
    \end{figure}
    \\



        \begin{figure}[H]
            \centering
            \begin{subfigure}[t]{0.49\textwidth}
                \centering
                \includegraphics[width=\linewidth, height=10cm, keepaspectratio=false]{9-zone-fausse.png}
                \caption{P09 : capture d'un artefact}
                \label{fig:artefact_p09}
            \end{subfigure}
            \hfill
            \begin{subfigure}[t]{0.49\textwidth}
                \centering
                \includegraphics[width=\linewidth, height=10cm, keepaspectratio=false]{carotid_mask_P09_after carotid _limit.png}
                \caption{P09 après limitation du masque}
                \label{fig:limitation_p09}
            \end{subfigure}
            \caption{Exemple de réduction des faux positifs (artefact éliminé) : comparaison des masques avant et après limitation pour le patient P09.}
            \label{fig:comparaison_p09}
        \end{figure}


        \subsubsection*{Défis et Limites}
            \textbf{Problèmes de clustering}\\
            L'approche MeanShift présente des limitations intrinsèques :
            \begin{itemize}
                \item Sensibilité au paramètre de bandwidth : des choix sous-optimaux peuvent fusionner artères et veines
                \item Suppose une distribution bimodale des intensités
                \item Coût computationnel élevé pour les volumes haute résolution (atténué par le sous-échantillonnage, au risque de perdre petits vaisseaux)
                \item N'impose pas de contraintes de tubularité, pouvant générer des masques fragmentés
            \end{itemize}
            \bigskip

            \textbf{Problèmes généraux}\\
            Cette approche dépend fortement du contraste des scans :
            \begin{itemize}
                \item Performance dégradée pour les low-dose scans (faible contraste vasculaire)
                \item Exemple visible avec le patient P03 : segmentation carotidienne obtenue mais intensités plus faibles que les données manuelles.
                Image de fond : moyenne d’intensité sur l’ensemble des volumes temporels (4D). L’arrière-plan est sombre en raison du faible contraste (low-dose scan) (cf Figure 16). La segmentation carotidienne est affichée en viridis, et les points les plus lumineux (en blanc) indiquent la position du cathéter
            \end{itemize}
            \\
            
            \begin{figure}[H]
                \centering
                \includegraphics[width=0.8\linewidth]{carotid_mask_P03.png}
                \caption{Exemple de low-dose scan P03. .}
                \label{fig:lowdose_p03}
            \end{figure}


        \subsubsection*{Conclusion}
            Bien que cette méthode donne des résultats satisfaisants pour les scans standard avec anatomie vasculaire typique (cf Figure 17), ses limitations (sensibilité au contraste, fragmentation des masques) ont motivé l'exploration d'autres approches, notamment la \textit{Ward parcellation} afin d'améliorer la cohérence spatiale par le biais du clustering hiérarchique.\\
            
            
         \noindent
        \begin{figure}[H]
            \centering
            \includegraphics[width=\linewidth]{TAC_fit_P04.png}
            \caption{Exemple de la courbe TAC ajustée pour le patient P04.}
            \label{fig:tac_p04}
        \end{figure}


    
    \hspace{1.5cm}
    %\hrule

%%%-------------------------------------------WARD----------------------------------------------------%%%

\newpage
        
        \subsubsection{CLUSTERING WARD \href{run:Pistes_et_Tests/piste_ward/2clustersNIlearn.py}{(\underline{\texttt{2clustersNIlearn.py}}})}\\
            Pour trouver les clusters ou les régions les plus intéressantes, on applique d'abord un masque binaire\footnote{Masque binaire: permet de supprimer les valeurs inférieures à un seuil pour ne garder que les données significatives.} afin de réduire les zones d'intérêts qu'on donne ensuite à Parcellations (roi\_mask dans la Figure 18).

            La méthode Parcellation() est particulièrement pertinente car elle nous donne le choix du nombre de clusters à renvoyer (n\_parcels dans la Figure 18), et du type d'algorithme qui nous convient le plus, dans cette approche nous avons choisi "ward".

            \subsubsection*{Algorithme et idée de l'approche}
                \begin{figure}[H]
                \centering
                \includegraphics[width=0.5\linewidth]{Parcellations.png}
                \caption{exemple d'utilisation de la méthode prédéfinie Parcellation}
                \label{fig:enter-label}
            \end{figure}
            \vspace{0.2cm}
            
                \subsubsection*{Principe de l'algorithme de Ward}
                \vspace{0.2cm}
                L'algorithme de Ward est une approche ascendante hiérarchique qui suit une logique similaire à celle des algorithmes gloutons. Le but est donc d’obtenir un nombre détérminé de clusters choisi par nous, tout en minimisant la variance à l’intérieur de chaque cluster.
                \vspace{0.2cm}

          Pseudo-code de l'algorithme:
            \begin{itemize}
                \item Extraire les caractéristiques des voxels (coordonnées spatiales + intensité).
                \item Initialiser chaque voxel comme un cluster individuel.
                \item Répéter jusqu'à atteindre le nombre de clusters souhaité :
                \begin{itemize}
                    \item Calculer la distance entre tous les clusters (basée sur l’inertie intra-cluster).
                    \item Fusionner les deux clusters dont la fusion augmente le moins la variance totale.
                \end{itemize}
                \item Retourner les clusters finaux représentant des régions cérébrales cohérentes.
            \end{itemize}

        
            \subsubsection*{Pourquoi on a choisi Ward pour 'Parcellations' ?}
                La méthode de Ward (cf Figure 19) présente plusieurs avantages, elle permet d'obtenir des clusters compacts, c’est-à-dire des groupes de voxels spatialement proches, ce qui facilite leur interprétation. De plus, cette méthode tend à regrouper des régions où l’activité présente des similarités temporelles ou des points dont les caractéristiques sont cohérentes entre eux, renforçant ainsi la pertinence fonctionnelle des clusters obtenus. Enfin, elle respecte globalement la structure anatomique du cerveau, ce qui est essentiel dans notre cas pour identifier avec plus de précision des zones comme les carotides.

                \begin{figure}[h]
                        \centering
                        \includegraphics[width=0.8\linewidth]{imagecomp.png}
                        \caption{comparaison entre les méthodes possibles de Parcellations}
                        \label{fig:enter-label}
                    \end{figure}
                    \label{fig:enter-label}

   
    
            \subsubsection*{Implémentation }
    
            \subsubsection*{Pré-traitement}
                Pour que ça soit plus facile et moins couteux à Parcellations de traiter les données, nous avons implementé \textbf{le masque binaire} pour diminuer le nombre de points à traiter à ceux potentiellement pertinents avec un seuil déterminé.

                \vspace{0.2cm}

                De plus, comme la fonction \texttt{Parcellations} ne peut traiter que des données en 3D (sans dimension temporelle), nous avons \textbf{sélectionné le volume temporel} que nous estimons le plus pertinent pour la détection des carotides. Généralement, nous sélectionnons le premier volume présentant de fortes intensités dans la région inférieure du cerveau, en suivant l’ordre chronologique.
                Cette opération a été réalisée à l'aide de la fonction \texttt{index\_img} de Nilearn qui permet d'extraire les données à un instant( coordonnée temporelle $t$) choisi.

\vspace{0.2cm}

                    
            \subsubsection*{Après \texttt{Parcellations}:}
                Une fois que les clusters sont trouvés, on les classe par ordre décroissant d'intensité. Ensuite, on identifie le cluster avec la plus forte intensité, ce qui correspond à la localisation de la carotide.
                Afin d’assurer que le cluster choisi appartient à l'hémisphère cérébral concerné (gauche ou droit), nous avons limité notre recherche aux extrémités coronales (selon l’axe des abscisses) définies par le masque du cerveau\footnote{Masque du cerveau: outil fourni par la fonction \textbf{compute\_brain\_mask} de Nilearn, utilisé pour identifier automatiquement la région du cerveau dans une image.}, afin de corriger certaines détections erronées dues à la présence de cathéters mal interprétés comme des carotides chez certains patients.

    
            \subsubsection*{Résultats}
                Visualisations :
                On obtient les coordonnées 3D de toutes les points appartenant aux clusters trouvés et on l’affiche pour visualiser le résultat pour le TAC ( cf Figure 20, Figure 21 , Figure 22) . 
                    \begin{figure}[H]
                    \centering
                    \includegraphics[width=0.75\linewidth]{clusterss.png}
                    \caption{visualisation des clusters produits avec Parcellations }
                    \label{fig:enter-label}
                \end{figure}
                \begin{figure}[H]
                    \centering
                    \includegraphics[width=0.75\linewidth]{resultatscluster.png}
                    \caption{Exemple de cluster détecté dans l’hémisphère droit}
                    \label{fig:enter-label}
                \end{figure}
                 \begin{figure}[H]
                \centering
                \includegraphics[width=0.75\linewidth]{resulTACclus.png}
                \caption{Après avoir choisi le cluster de la carotide, on représente la TAC de l'intensité moyenne de tous les points appartenant à ce cluster.}
                \label{fig:enter-label}
                 \end{figure}

                
                
                
            \subsubsection*{Limites}
                Discussions des limitations :
                	 La méthode n'est pas entièrement automatique: il faut trouver manuellement le seuil en plus que le volume à utiliser pour les données 3D pour chaque patient.

            \subsubsection*{Conclusion}
                Pour fixer ces limitations, nous avons décidé de combiner plusieurs méthodes afin d'améliorer l'automatisation.
                
    \hspace{1.5cm}

%%%-------------------------------------------idif_pipeline_pet----------------------------------------------------%%%
\newpage
\subsection{Méthode finale}
    \emph{Fichier : \href{run:idif_pipeline_pet.py}{\underline{\texttt{idif\_pipeline\_pet.py}}}}

    \subsection*{Objectif général}
        \noindent
        L'objectif de ce pipeline est d'extraire automatiquement une courbe d'entrée vasculaire (IDIF) à partir d'images TEP dynamiques 4D, sans intervention manuelle. Contrairement aux autres essais basés sur \textit{K-Means} ou \textit{Mean Shift}, ce pipeline constitue une version améliorée de la méthode explorée précédemment. Il repose sur une segmentation par clustering hiérarchique (\textit{Ward}) et intègre une modélisation de type \textit{Feng} pour extraire des courbes temporelles physiologiquement exploitables.
 

    \subsection*{L'importance du double masquage}

        \noindent
        Le prétraitement commence par la création d’une image de référence, obtenue en moyennant des frames temporelles. Cette opération permet de lisser les fluctuations superflues, tout en mettant en évidence les structures fixes du cerveau. L’image résultante est plus stable et sert de base pour la segmentation.\\

        \noindent
        À partir de cette image moyennée, un \textbf{masque cérébral} \textsuperscript{figure 23} est généré automatiquement à l’aide de la fonction \texttt{compute\_brain\_mask} fournie par \textit{Nilearn}. Un masque est une image binaire (valeurs 0 ou 1) qui permet d’isoler les voxels appartenant à une région d’intérêt. Dans ce cas, le masque identifie les voxels du volume cérébral, en se basant sur les niveaux d’intensité.
            \begin{figure}[H]
                    \centering
                    \includegraphics[width=0.75\linewidth]{brain_mask.png}
                    \caption{Brain mask crée}
                    \label{fig:enter-label}
                \end{figure}

        \noindent
        Cependant, ce masque brut peut contenir du bruit ou des discontinuités. C’est pourquoi il est post-traité à l’aide d’opérations morphologiques classiques : une \textit{fermeture} (closing) pour combler les petits trous internes, et une \textit{érosion} pour éliminer les artefacts isolés ou les bords irréguliers. Ces étapes garantissent un masque plus cohérent sur le plan anatomique.\\

        \noindent
        Ensuite, une seconde restriction est appliquée : on extrait uniquement la \textbf{partie inférieure du cerveau} (zone basale) à l’aide de la fonction \texttt{get\_lower\_brain\_mask} (cf Figure 24). Cette étape repose sur une observation anatomique : les segments intracrâniens des artères carotides internes sont situés à la base du cerveau, en profondeur. En limitant notre attention à cette zone, on réduit le bruit potentiel dû à d'autres structures brillantes (comme le cortex ou les ventricules), et on améliore la pertinence du clustering à venir.
                    \begin{figure}[H]
                    \centering
                    \includegraphics[width=0.75\linewidth]{partie inférieur du cerveau.png}
                    \caption{Brain mask crée}
                    \label{fig:enter-label}
                \end{figure}

        \noindent
        En résumé, cette double étape de masquage permet de restreindre l’analyse à une région anatomiquement pertinente, tout en éliminant les artefacts indésirables et les zones peu informatives de l’image TEP.


    \subsection*{Sélection automatique d’une frame pertinente}
        
        \noindent
        Pour identifier au mieux les structures vasculaires, il est essentiel de choisir une frame temporelle où les carotides sont bien visibles. Le choix de cette frame repose sur une analyse de l’intensité moyenne dans la zone inférieure du cerveau. En effet, les artères intracrâniennes présentent une activité très élevée aux premiers instants de l’injection du traceur.
        
        \noindent
        Le pipeline parcourt ainsi toutes les frames et calcule leur intensité moyenne au sein de la région basale. La première frame dépassant 70\% de l’intensité maximale est retenue. Cette approche simple et robuste permet d’automatiser la sélection d’une image TEP pertinente pour la suite du traitement.
        
    \vspace{0.5em}
    \subsection*{Détection des régions actives par seuillage}
        
        \noindent
        À partir de la frame sélectionnée, on applique un seuillage basé sur un percentile élevé (96.9\%) pour isoler les voxels les plus brillants, qui sont supposés correspondre à des structures fortement vascularisées. Cette étape permet de réduire drastiquement le nombre de voxels analysés tout en conservant les zones les plus informatives.
        
        \noindent
        Le masque binaire ainsi obtenu met en évidence des régions d’activité intense, constituant une bonne base pour une parcellisation ultérieure.
        
        \vspace{0.5em}
        \subsection*{Parcellisation hiérarchique et sélection anatomique}
        
        \noindent
        Sur ce masque, une parcellisation est effectuée via le clustering hiérarchique de type \texttt{Ward}. Cette méthode regroupe les voxels proches spatialement et présentant une intensité homogène, en minimisant la variance intra-cluster. (ici c'est répétition)
        
        \noindent
        Chaque cluster est ensuite évalué selon deux critères :
        \begin{itemize}
            \item Son intensité moyenne sur l’image 3D.
            \item Sa position spatiale dans le cerveau (hémisphère gauche ou droit).
        \end{itemize}
        
        \noindent
        On retient alors les deux clusters les plus brillants, respectivement dans chaque hémisphère, en s’assurant qu’ils se situent bien dans les zones attendues des carotides internes.
        
        \vspace{0.5em}
    \subsection*{Extraction et modélisation des courbes TAC}
        
        \noindent
        À partir des voxels de chaque cluster carotidien (gauche, droit, combiné), on extrait une courbe TAC (Time-Activity Curve), représentant l’évolution temporelle de la concentration du traceur. Ces courbes sont essentielles pour les étapes ultérieures de quantification TEP.
        
        \noindent
        Pour stabiliser les TAC et les rendre plus exploitables, un ajustement est effectué à l’aide du modèle gamma-varié de Feng:
        
        \[C(t) = (A_1 - A_2 - A_3) \cdot e^{-\lambda_1 t} + A_2 \cdot e^{-\lambda_2 t} + A_3 \cdot e^{-\lambda_3 t}\]

        
        \noindent
        Celui-ci est appliqué uniquement à la phase décroissante de la courbe — après le pic d’intensité — afin d’éviter les erreurs liées à la montée initiale souvent bruitée. La phase ascendante est remplacée par des zéros pour garantir une continuité temporelle sans surinterprétation.
        
    \vspace{0.5em}
    \subsection*{Visualisation interactive et export}
        
        \noindent
        Les résultats sont résumés dans une figure complète (cf Figure 25) générée par la fonction \texttt{plot\_results}, qui présente :
        \begin{itemize}
            \item Une visualisation 3D des clusters identifiés comme carotides.
            \item Les courbes TAC expérimentales et leurs ajustements pour chaque côté et pour les deux carotides combinées.
            \item Une comparaison facultative avec une courbe de référence manuelle (.crv), activable par la touche '\texttt{a}'.
        \end{itemize}
        
        (mettre une image ex)
    \begin{figure}[H]
                \centering
                \includegraphics[width=0.9\linewidth]{TACs_p01.png}
                \caption{Visualisation sur p01}
                \label{fig:enter-label}
            \end{figure}
        
        \vspace{0.5em}
    \subsection*{Conclusion}

        \noindent
            Cette méthode automatique basée sur le clustering Ward et le fit de Feng nous a permis de construire un pipeline complet et fonctionnel pour extraire une IDIF sans intervention manuelle. Elle représente une version aboutie de notre travail après plusieurs essais et ajustements.
        
        \noindent
            Elle sera utile à présenter pour montrer l’évolution de notre réflexion, les différentes étapes testées, et comment on a utilisé les bons outils pour arriver à un résultat cohérent. La suite du rapport présente justement ces outils en détail.

        
                
            
        

%%%-------------------------------------------Outils----------------------------------------------------%%%
    \newpage
    \section{Outils et environnement de développement}
    \label{sec:Outils}
    \noindent
        Pour mener à bien ce projet, nous avons utilisé différents outils pour le traitement des images TEP, le développement du pipeline et le travail en équipe.\\
    
        \noindent
        Les données d’entrée sont des fichiers NIfTI (\texttt{.nii}) contenant les images TEP dynamiques 4D de plusieurs patients. Pour les visualiser, nous avons utilisé le logiciel open-source \textbf{ITK-SNAP}, recommandé par notre encadrant. Ce logiciel permet d’explorer facilement les volumes médicaux en naviguant dans les plans axiaux, sagittaux et coronaux, ainsi que dans les différentes frames temporelles. Il s’est révélé particulièrement utile pour localiser les artères carotidennes.\\
        
        \noindent
        L’ensemble du pipeline a été codé en \textbf{Python}, un langage bien adapté à ce type de traitement grâce à sa simplicité, sa flexibilité et son écosystème riche en bibliothèques spécialisées. Parmi celles-ci :
    \begin{itemize}
        \item \textbf{Nilearn} a été utilisée pour manipuler et visualiser les images NIfTI. Elle nous a également permis d’explorer différentes méthodes et algorithmes en lien avec notre objectif final, grâce à ses nombreuses fonctions prédéfinies adaptées à la segmentation et au traitement des images cérébrales.
        \item \textbf{Scipy.ndimage} nous a permis d’appliquer des opérations morphologiques comme l’érosion ou la fermeture.
        \item \textbf{Skimage} a fourni des outils puissants pour la segmentation et le traitement des régions.
        \item \textbf{Matplotlib} a servi à tracer les courbes d’intensité (TAC) et visualiser les résultats.
        \item \textbf{Scikit-learn} a été utilisée pour appliquer les algorithmes de clustering \texttt{KMeans}, \texttt{Mean Shift} et \texttt{Ward}.\\
    \end{itemize}
    
        \noindent
        Un problème rencontré concerne la gestion des fichiers volumineux : Git ne prend pas en charge les fichiers \texttt{.nii} et \texttt{.crv}, ce qui compliquait leur partage. Pour résoudre cela, nous avons développé un script \href{run:./download_data.py}{\textbf{\texttt{download\_data.py}}} permettant de télécharger automatiquement les données des patients depuis un dossier partagé sur Google Drive. Ce script doit être lancé avant d’exécuter le pipeline principal \href{run:./idif_pipeline_pet.py}{\texttt{idif\_pipeline\_pet.py}}.\\

    
        \noindent
        Enfin, pour assurer une bonne collaboration tout au long du projet, nous avons utilisé \textbf{GitHub} (via GitHub Desktop) pour la gestion de version, ce qui nous a permis de suivre l’évolution du code de manière centralisée et de travailler efficacement en équipe. La communication avec notre encadrant s’est déroulée via des visioconférences régulières sur \textbf{Discord}, facilitant les points d’avancement fréquents et les retours rapides.\\
        
         \noindent
        Nous avons également rédigé ce rapport de manière collaborative sur Overleaf, via la plateforme web \\ \href{https://plmlatex.math.cnrs.fr/}{\texttt{https://plmlatex.math.cnrs.fr/}}, qui nous a permis d’éditer le document à plusieurs en temps réel.\\


        \noindent
        Enfin, nous prévoyons d’utiliser Canva pour préparer les slides de la présentation finale, afin de garantir une mise en page claire, attractive et professionnelle.\\
        
        \noindent
        Nous avons également fait appel à des outils d’intelligence artificielle, notamment \textbf{Mistral} ou \textbf{ChatGPT}, pour nous aider à comprendre certaines parties complexes du code, à corriger des erreurs, à commenter des fonctions, ou encore à structurer notre travail. Cette aide nous a permis de gagner du temps, mais aussi de mieux assimiler les concepts techniques au cœur du projet.

%%%-------------------------------------------Analyse----------------------------------------------------%%%

\newpage
\section{Analyse des résultats}\\
\noindent
Exemple de comparaison des courbes résultantes des différentes pistes sur les p01 et p04 (cf Figure 26 et Figure 27):\\

    \begin{figure}[H]
        \centering
        \includegraphics[width=0.75\linewidth]{figure01analyse.png}
        \caption{Comparaison sur le patient p04}
        \label{fig:enter-label}
    \end{figure}
    
    \begin{figure}[H]
        \centering
        \includegraphics[width=0.75\linewidth]{figure2analyse.png}
        \caption{Comparaison sur le patient p01}
        \label{fig:enter-label}
    \end{figure}
        
        \vspace{0.5em}
        
        

        On remarque que Kmeans obtient des valeurs bien supérieures aux autres et c’est normal, dans cette version nous avons choisi de prendre la valeur la plus élevé du cluster afin d’avoir un résultat plus clair et pas un point excentré du centre du cluster qui serait peut être plus insignifiant.\\
        
        \noindent
        En comparaison, \texttt{MeanShift} et \href{run:./idif_pipeline_pet.py}{\texttt{idif\_pipeline\_pet.py}} ont des résultats beaucoup plus proches et peut-être plus réalistes.\\
        Pour les résultats des analyses, nous allons donc nous concentrer sur \href{run:./idif_pipeline_pet.py}{\texttt{idif\_pipeline\_pet.py}}, qui est notre fonction finale, les deux autres n'étant que des versions expérimentales.\\
        
        \subsection*{\href{run:./idif_pipeline_pet.py}{\texttt{idif\_pipeline\_pet.py}} :}\\

        
        \begin{figure}[H]
            \centering
            \includegraphics[width=0.75\linewidth]{figure3analyse.png}
            \caption{Comparaison des différentes approches}
            \label{fig:enter-label}
        \end{figure}
        
        \noindent
        Dans la Figure 28 pour le patient p05, on trouve les 2 carotides ( droite et gauche )

        
        \begin{figure}[H]
            \centering
            \includegraphics[width=0.75\linewidth]{figure4analyse.png}
            \caption{Cluster identifiés}
            \label{fig:enter-label}
        \end{figure}
        
        \noindent
        La Figure 29 montre les différents clusters identifiés par la méthode, on voit qu’on obtient des clusters proches mais des clusters qui semblent être positionnées sur la carotide droite et gauche, on a donc les positions des 2 carotides et on peut en tirer les résultats suivants :

        
        \begin{figure}[H]
            \centering
            \includegraphics[width=0.8\linewidth]{figure5analyse.png}
            \caption{Exemple de sortie de notre méthode finale pour le patient 04}
            \label{fig:enter-label}
        \end{figure}
        
        \noindent
        Grâce à cette méthode on peut extraire la TAC de chaque carotide et produire une TAC qui est la moyenne des 2 carotides (cf Figure 30).\\
        On peut voir que les valeurs de la carotide droite augmente fortement puis redescendent jusqu’a 0, tandis que la carotide gauche a un pic plus léger mais conserve une concentration dans sa carotide gauche.\\
        
        \noindent
        La différence observée entre les deux courbes de TAC pourrait s’expliquer par plusieurs facteurs  : des différences anatomiques ( origine aortique directe pour la carotide gauche, origine via le tronc brachiocéphalique pour la carotide droite), mais surtout des différences hémodynamiques et perfusion locales, et potentiellement des effets techniques liés à l’acquisition.\\
        
        \noindent
        Le pic plus élevé à droite (cf Figure 31) suivi d’une chute rapide peut indiquer un artefact de mesure, une perfusion plus rapide, ou un volume sanguin plus faible à cet endroit, tandis que la gauche montre un plateau plus prolongé, suggérant une clairance plus lente ou un volume plus important.\\
        
        \begin{figure}[H]
            \centering
            \includegraphics[width=0.8\linewidth]{p10analyse.png}
            \caption{Exemple de sortie avec comparaison pour le patient 10}
            \label{fig:enter-label}
        \end{figure}
        
      \noindent
      On constate que les valeurs obtenues pour la carotide droite sont globalement similaires à celles fournies par la TAC de référence. Toutefois, selon les patients, notre programme final peut donner des valeurs légèrement plus ou moins élevées. Malgré ces écarts, les courbes présentent une même tendance.\\
        
        


    

%%%-------------------------------------------Gestion----------------------------------------------------%%%
\newpage

\section{Gestion du Projet}

    \noindent
    Dès le lancement du projet, nous avons structuré notre travail en plusieurs phases successives.
    
    \begin{enumerate}
        \item Un premier sprint en décembre a été consacré à la prise en main des données TEP, à l’exploration des différentes approches de segmentation, à l’utilisation du logiciel ITK-SNAP, ainsi qu’à quelques tests préliminaires. Cette phase s’est conclue par un rapport de sprint (voir \hyperref[bibliographie]{\textbf{la bibliographie}}).
        \item La deuxième phase, menée entre février et mars, s’est concentrée sur l’implémentation et la comparaison des différentes méthodes de segmentation (KMeans, Mean Shift, Ward). Nous avons présenté nos avancées sous forme de diaporama conçu avec Canva, lors d’un échange avec notre encadrant (également voir \hyperref[bibliographie]{\textbf{la bibliographie}}).
        \item Enfin, la phase finale, entre avril et début mai, a été dédiée au développement du pipeline final, à sa validation sur plusieurs patients, à l'analyse des résultats ainsi qu’à la rédaction du rapport et l'organisation finale de la structure de notre projet.
    \end{enumerate}
    
    \noindent
    Les tâches ont été réparties entre les membres du groupe, et un suivi rigoureux a été mis en place grâce à Git. Tous les membres sont restés investis depuis le début du projet. Des réunions hebdomadaires ont été organisées pour coordonner le travail, partager les avancées et résoudre collectivement les blocages techniques. À chaque fin d’étape, des rendez-vous en ligne avec notre encadrant nous ont permis d’obtenir des retours structurés et de réajuster nos priorités si nécessaire.\\
    
    \noindent
    Pour faciliter la collaboration, nous avons centralisé notre code sur GitHub via l’interface GitHub Desktop. Cela nous a permis d’assurer un versioning clair, de prévenir les conflits, et de garantir un accès commun aux différentes versions du projet, en assurant une continuité de développement entre les différentes phases.\\
    
    \noindent
    Plusieurs choix initiaux ont été ajustés au fil du projet. Par exemple, bien que nous ayons d’abord testé les algorithmes KMeans et Mean Shift, ces derniers se sont avérés peu fiables dans notre contexte. Nous avons donc opté pour la parcellisation hiérarchique de Ward, plus cohérente et robuste pour identifier les régions d’intérêt vasculaire dans les images TEP.\\
    
    \noindent
    Avec davantage de temps, nous aurions pu aller plus loin dans la validation du code en développant une suite de tests unitaires, afin de renforcer la fiabilité de chaque fonction du pipeline. Cela aurait également permis de détecter plus rapidement les effets de modifications futures.\\
    
    \noindent
    Par ailleurs, pour améliorer la modularité de notre projet, il aurait été pertinent de répartir les différentes fonctions du fichier principal dans plusieurs fichiers Python dédiés (par exemple : segmentation.py, visualisation.py, fit-model.py, etc.). Cette organisation aurait facilité les modifications ciblées, la réutilisabilité du code, et l’intégration dans un environnement plus vaste.\\
    
    \noindent
    Enfin, une perspective importante pour la suite serait l’intégration du pipeline dans le logiciel interne du CHU de Montpellier, afin de faciliter l’analyse des données TEP dynamiques dans le cadre de la recherche clinique. Ce serait une manière concrète de valoriser le travail réalisé en le mettant au service de la communauté médicale.
      
%%%-------------------------------------------Conclusion----------------------------------------------------%%%
\newpage
\section{Bilan et Conclusions}
    Ce projet a constitué une opportunité enrichissante d’appliquer nos compétences en informatique à un domaine concret et exigeant : l’imagerie médicale. L’objectif principal, à savoir l’extraction automatique de courbes temps-activité (TAC) à partir de fichiers TEP dynamiques, a été atteint à travers une exploration méthodique de différentes approches de traitement et d’analyse d’image. \\
    
    \noindent
    Nous avons su nous adapter aux spécificités des données médicales en testant plusieurs méthodes de segmentation et de clustering, avant de converger vers un algorithme heuristique optimal. \\
    
    \noindent
    Ce projet nous a permis de nous familiariser avec des formats de données et logiciels scientifiques nouveaux, ainsi qu’avec des méthodes avancées de traitement d’image et d’apprentissage non supervisé. Au-delà de l’aspect technique, il a également mobilisé des compétences transversales importantes : gestion de projet, documentation des avancements, vulgarisation scientifique, collaboration en équipe et communication avec un partenaire extérieur à l’Université de Montpellier. \\
    
    \noindent
    Les principales difficultés rencontrées ont été liées à la complexité du contexte biomédical et à la diversité des cas patients. Néanmoins, ces obstacles ont été surmontés par un travail collectif efficace et un dialogue constructif avec notre encadrant. \\
    
    \noindent
    Plusieurs perspectives d'amélioration peuvent être envisagées. Par exemple, le projet s’est appuyé principalement sur une validation visuelle des TAC. Une amélioration notable consisterait à développer des métriques quantitatives (corrélation, aire sous la courbe, erreur quadratique moyenne, etc.) pour comparer les courbes extraites automatiquement à des références manuelles. \\
    De plus, le projet s’est concentré sur une zone d’étude réduite. Nous pourrions étendre les recherches à d’autres zones d'intérêt que la carotide mais surtout sur un échantillon plus vaste et hétérogène de patients afin de confirmer la pertinence de la méthode. \\
    Enfin, une interface graphique intuitive à destination des non-informaticiens rendrait ce projet utile et accessible. \\
    
    \noindent
    En somme, ce projet a été formateur tant sur le plan technique que méthodologique. Il nous a permis de mieux appréhender les exigences d’un projet interdisciplinaire, en nous confrontant à des problématiques réelles où l’informatique est mise au service de l’innovation médicale. \\




%%%-------------------------------------------Bibliographie----------------------------------------------------%%%
\newpage
\begin{thebibliography}{9}
\section*{Bibliographie}
\label{bibliographie}
  
\bibitem{Article}
\bibitemAticle{Aticle about \textit{"Image-derived input functions for human PET studies"} (2023)}  \url{https://pubmed.ncbi.nlm.nih.gov/37947880/}

\bibitem{meanshift}
Region Growing and Clustering Segmentation. \url{https://sbme-tutorials.github.io/2019/cv/notes/6_week6.html#mean-shift-clustering}

\subsection*{Annexes : }
\textit{Rapport d'avancement 1} : \href{../Rapports/Rapport_Sprint_Décembre.pdf}{CR Sprint Décembre}\\
\textit{Rapport d'avancement 2} : \href{../Rapports/Rapport_avancement_post_sprint2.pdf}{CR Post Sprint 2}\\
\end{thebibliography}

    
\end{document}
